\subsection{A bounded audit of raw multiplication integrality}

This \fid{F65} supplement records an unsuccessful
search for a fractional-multiplicity or non-Laurent
obstruction among the actual printed raw constants.
It is not a general integrality theorem, and adds
no such source counterexample to the main argument.
Conjecture~3.4, \S3.5, p.~24 of VIII uses the
positive-factor convention defined near
Conjecture~2.7, \S2.4.2, p.~15: positive Laurent
coefficients are permitted to be rational
\cite{VIII}. An actual finite Euler realization
would impose the additional grouped-integrality
condition proved below.

\paragraph{The six raw formulas actually printed.}
VIII, \S5.1.2, p.~29 defines
\(c(b,b';b'')\) by expansion of \(bb'\).
Its Figure~2, p.~34, contains thirteen labeled
entries covered by these six distinct formulas
\cite{VIII}. The audit letters in this table
are not replacements for the source's basis labels.
The first row's printed gamma-label irregularity
need not be resolved to read its formula.
\begin{center}
\begin{tabular}{c l c}
\toprule
Audit name & Raw formula in \(q\) & Number of labels\\
\midrule
\(a\)&\(- (1+q)^2(q^2+q+1)(q-1)^2/q^3\)&3\\[3pt]
\(b\)&\(- (q^2+q+1)(q-1)^2/q^2\)&1\\[3pt]
\(c\)&\((1+q)^2/q\)&4\\[3pt]
\(d\)&\((1+q)^2(6q+5q^2+5)/q^2\)&2\\[3pt]
\(e\)&\((1+q^2)^2(q^2+q+1)(q-1)^2/q^4\)&1\\[3pt]
\(f\)&\((1+q)^2(2q^2+q+1)(q^2+q+2)/q^3\)&2\\
\bottomrule
\end{tabular}
\end{center}
All six belong to \(\Z[q,q^{-1}]\).
Their full coefficient vectors, in increasing
exponent order and retaining interior zeros, are
\[
\begin{array}{ccl}
a:&[-3,3]&(-1,-1,1,2,1,-1,-1),\\
b:&[-2,2]&(-1,1,0,1,-1),\\
c:&[-1,1]&(1,2,1),\\
d:&[-2,2]&(5,16,22,16,5),\\
e:&[-4,4]&(1,-1,2,-3,2,-3,2,-1,1),\\
f:&[-3,3]&(2,7,14,18,14,7,2).
\end{array}
\]
The displayed powers of \(q\) in the denominators
are Laurent exponents, not fractional multiplicities.

Nor does exact removal of the prescribed
half-step vanishing factor create fractions in
these examples. Put \(u=q^{1/2}\).
If an integral Laurent polynomial equals
\((u-u^{-1})^s\) times a rational Laurent
polynomial \(g/N\), choose the least positive
common coefficient denominator \(N\), so
\(g\in\Z[u,u^{-1}]\).
If a prime \(p\) divides \(N\), reducing the
identity modulo \(p\) gives
\((u-u^{-1})^s\bar g=0\).
The first factor is nonzero in the integral
domain \(\F_p[u,u^{-1}]\), so \(\bar g=0\).
Every coefficient of \(g\) would be divisible
by \(p\), contradicting the minimality of \(N\).
Thus \(N=1\), and the quotient is integral.
This concerns an exact Laurent quotient, not
arbitrary rational-function normalization.

\paragraph{Why the other displayed denominators do not answer this test.}
The bounded audit also checked all twelve
tables in Figures~24--35, pp.~58--65,
containing 542 signed integer entries.
These are normalized vectors, sometimes of
field traces or norms, not a raw unnormalized
multiplication table.
In \S5.4.1, pp.~48--49, the constants are
rescaled by powers of \(\hat f_p(q)\), with
individual exponents not supplied.
Section~5's vector convention, p.~27, then
strips a maximal vanishing factor and records
only half of the resulting bar-invariant vector.
These data do not prove Laurentness or integrality
of the omitted raw constants.

Section~5.2 begins on p.~29; its
Figures~3--5, pp.~35--37, are ambient-expansion
vectors for one element at \(r=4\).
The explicit halves and algebraic coefficients
in Figure~6, p.~38, belong to the ambient
monomial-basis expansion discussed in
\S5.3, p.~31.
The nonmonomial denominators on p.~33 occur
in a balance-fiber change of basis.
None of these entries is thereby a raw
\(c(b,b';b'')\) multiplication constant
\cite{VIII}.
No unprinted constant or omitted rescaling
exponent was recovered by this audit.

\paragraph{The precise grouped-integrality test.}
Let \(f(z)=\sum_j a_jz^j\in\Q[z,z^{-1}]\)
and fix a nonzero algebraic \(\beta\).
Suppose for every sufficiently large integer \(n\)
that \(f(\beta^n)\) is the uninserted Euler
trace of one bounded finite-dimensional point
complex with invertible Frobenius.
For each distinct root \(\gamma\), one must have
\begin{equation}\label{eq:VIIIGroupedIntegerMultiplicity}
       \sum_{j:\,\beta^j=\gamma}a_j\ \in\ \Z.
\end{equation}
This assertion needs neither purity nor
Frobenius semisimplicity.

\begin{proof}
The Euler trace of the complex is a finite
power sum whose grouped coefficients are
differences of actual eigenvalue multiplicities,
hence integers. Subtract it from
\(\sum_j a_j(\beta^j)^n\).
For the distinct nonzero roots in their union,
the matrix of sufficiently large consecutive
powers is a Vandermonde matrix times an
invertible diagonal matrix. It is invertible,
so every grouped coefficient agrees.
This proves \eqref{eq:VIIIGroupedIntegerMultiplicity}.
\end{proof}

If \(\beta\) is not a root of unity, its
integer powers are distinct and every \(a_j\)
must be integral. If it has finite order \(h\),
only the residue-class sums modulo \(h\)
are forced to be integral.
Pointwise integer values alone are insufficient:
\((1+(-1)^n)/2\) is integer-valued but has
coefficient \(1/2\) at each of two distinct
roots, so it is not such an uninserted Euler trace.
An overall Tate or cohomological shift changes
the roots or an overall sign, without fixing
a fractional grouped multiplicity.
Inserted-operator traces need not have integral
coefficients and are outside this test.
For a genuinely non-Laurent rational function
at a fixed real dilation \(\beta>1\),
\fid{F26} gives a separate finite-spectrum
obstruction; none was found among the raw
formulas checked here.

\paragraph{Actual mixed scalar witnesses for all six formulas.}
Let \(Q\) be a prime power, \(\ell\nmid Q\),
and \(E/\Q_\ell\) finite. Under the explicitly
chosen law \(q_{\rm src}=Q^n\), every one
of the six formulas has a bounded mixed
uninserted Tate realization.
For \(f(q)=\sum_j a_jq^j\) define
\[
 K_f=
 \bigoplus_{a_j>0}E(-j)^{\oplus a_j}
 \ \oplus\!
 \bigoplus_{a_j<0}E(-j)^{\oplus(-a_j)}[-1].
\]
The first summands have degree zero, the
second degree one. Geometric Frobenius on
\(E(-j)\) is \(Q^j\), so
\[
       \operatorname{Tr}_{\rm Eul}(F^n;K_f)
                    =\sum_j a_jQ^{jn}=f(Q^n).
\]
These are actual finite ordinary Tate complexes,
not merely rationalized Grothendieck classes.
They do not identify multiplication maps,
designated source bases, canonical simples,
or one common pure weight.
Thus this bounded integrality audit supplies
no further obstruction to mixed scalar
realizations of the printed Figure~2 formulas.
Rescaling a designated basis to clear a
different denominator would change the
constants and require a new compatible
multiplication/action comparison.

% Separate paragraph for Part III.
\paragraph{The \fid{F65} bounded integrality audit in Grothendieck's language.}
Actual finite graded coefficient spaces have
integral dimensions. Rationalizing their
Grothendieck group permits rational virtual
classes without furnishing corresponding
finite objects.
SGA~6, V, \S\S2--3 and cote~133, scan~18
supply this coefficient representation language
\cite{SGA6,Archive}; no \(\lambda\)-ring structure
on the source noncommutative algebra is asserted.
The explicit mixed Tate witnesses belong to
SGA~6, I, Definition~4.7, pp.~121--122 and
Example~4.8, p.~122 \cite{SGA6}.
The source's permitted rational positive factor
alone supplies no Euler realization in that
framework. The test and bounded negative-search
scope remain distinct from the stronger pure
and operation comparisons elsewhere in \fid{F65}.
